\chapter{Deployment: Merging, Inference Memory and Latency}
\label{ch:deploy}

The previous chapters were about training. This chapter asks what LoRA changes once the model is trained and makes
predictions (inference). A LoRA model can be deployed in two forms: \emph{merged}, with the adapters added into the
weights, or \emph{unmerged}, with the adapters kept as separate layers. All measurements use the adapters trained
with seed 42; the procedures are described in Section~\ref{sec:procedures}.

\section{Merging leaves the predictions unchanged}

For all seven LoRA configurations of each model (three targets and four placements), merging does not change the
predictions:
\begin{itemize}
  \item the test accuracy is identical with the adapters unmerged, merged, and after 100 cycles of merging and
        unmerging, and every one of the 10,000 test images receives the same predicted class;
  \item on eight random inputs, the largest difference in any of the 100 class scores between the merged and the
        unmerged model is at most $5.0\times10^{-6}$ on ViT-B and $1.7\times10^{-6}$ on Swin-B, the size of
        floating-point rounding;
  \item 100 cycles of merging and unmerging change no weight by more than $1.2\times10^{-7}$.
\end{itemize}
The outputs are equal up to rounding, not bit for bit. The check uses one trained adapter per configuration;
because merging is exact algebra (Chapter~\ref{ch:background}), we do not expect other seeds to behave
differently. Table~\ref{tab:merge} in the appendix gives the values per configuration.

\section{Inference memory}

\begin{table}[!htb]
\centering\small
\begin{tabular}{lcccc}
\toprule
 & \multicolumn{2}{c}{ViT-B} & \multicolumn{2}{c}{Swin-B} \\
configuration & batch 1 & batch 64 & batch 1 & batch 64 \\
\midrule
\rawinput{data/tab_infer.tex}
\bottomrule
\end{tabular}
\caption{Peak inference memory (GiB) of a forward pass without gradients, at batch sizes 1 and 64 (profiler). For
the LoRA configurations: unmerged / merged. The other three placements give the same values as the odd blocks to
within 0.001~GiB.}
\label{tab:infer}
\end{table}

Inference memory hardly depends on LoRA (Table~\ref{tab:infer}).
\begin{itemize}
  \item A \textbf{merged} model needs exactly the memory of the dense model: 0.80~GiB (ViT-B) and 1.80~GiB (Swin-B)
        at batch size 64.
  \item \textbf{Unmerged attention adapters} add 1--3~MiB, the size of the adapters themselves.
  \item \textbf{Unmerged MLP adapters} add 0.15~GiB on ViT-B and 0.39~GiB on Swin-B at batch size 64, and about
        5--10~MiB at batch size 1. A likely reason, which we did not isolate, is that the adapter's output must exist
        next to that of the frozen layer before the two are added, and for the first MLP layer this output is four
        times wider than the block.
\end{itemize}
Inference memory is also far smaller than training memory (0.8 against 5.5--8.2~GiB on ViT-B at batch size 64),
because nothing is kept for a backward pass. LoRA therefore does \emph{not} lower the memory needed for
predictions. Its memory saving is a saving in training.

\section{Latency}
\label{sec:latency}

Every configuration was timed on the CIFAR-100 test images in four ways: one image at a time end to end, which is
the situation of a device that classifies images as they arrive; in batches of 1 to 128 images (throughput); with
the forward pass replayed as a \emph{CUDA graph}, a diagnostic explained below; and on the CPU.
Figure~\ref{fig:latency} shows the end-to-end time per image; Table~\ref{tab:latency} in the appendix gives all values.


\begin{figure}[!htb]
\centering
% generated by scripts/make_figures.py; do not edit
\begin{tikzpicture}
\begin{groupplot}[group style={group size=2 by 1, horizontal sep=\figSep, vertical sep=1.3cm}, width=\figW, height=\figH, repaxis, ybar, /pgf/bar width=9pt, ymin=0, xmin=0.4, xmax=4.6, xtick={1,2,3,4}, xticklabels={Attn,MLP,Attn+MLP,odd}, ymajorgrids, legend style={font=\footnotesize, draw=none, fill=none}, legend pos=north west, legend columns=-1, enlarge y limits={upper, value=0.25}]
\nextgroupplot[title={ViT-B}, ylabel={time per image, end to end (ms)}]
\addplot[fill=cGrey!35, draw=black] coordinates {(1,7.14) (2,7.14) (3,7.15) (4,7.13)};
\addlegendentry{merged}
\addplot[fill=cAccent, draw=black] coordinates {(1,8.29) (2,8.13) (3,9.23) (4,7.71)};
\addlegendentry{unmerged}
\draw[black, densely dashed, thick] (axis cs:0.4,7.12) -- (axis cs:4.6,7.12);
\addlegendimage{black, densely dashed, thick, line legend}
\addlegendentry{dense}
\nextgroupplot[title={Swin-B}]
\addplot[fill=cGrey!35, draw=black] coordinates {(1,9.80) (2,9.92) (3,9.83) (4,9.83)};
\addplot[fill=cAccent, draw=black] coordinates {(1,12.79) (2,12.14) (3,15.06) (4,11.12)};
\draw[black, densely dashed, thick] (axis cs:0.4,9.74) -- (axis cs:4.6,9.74);
\end{groupplot}
\end{tikzpicture}

\caption{Median time to classify one test image end to end on the GPU (9,950 images, mean of three processes).
``odd'' = attention LoRA on the odd blocks.}
\label{fig:latency}
\end{figure}

\paragraph{A merged model is as fast as a dense one.} End to end, merged models take 7.13--7.15~ms per image
against 7.12~ms for the dense ViT-B, and 9.75--9.92~ms against 9.74~ms for the dense Swin-B; the throughput at
batch size 64 is the same (305--306 and 262--263 images per second). On Swin-B the forward pass of the merged
models was 0.3--2\% slower than that of the dense model. This comes from the order of measurement: in two of the
three processes the merged model was timed after the unmerged one, and in the process that timed it first it is as
fast as the dense model. With the CUDA graph the two are equal (7.50--7.51~ms). Merging has no measurable cost at
prediction time.

\paragraph{Unmerged adapters slow down single-image predictions.} With the adapters kept separate, one image takes
7.7--9.2~ms on ViT-B ($+$8 to $+$30\% against dense) and 11.1--15.1~ms on Swin-B ($+$14 to $+$55\%). The slowest 1\%
of images follow the same pattern. The penalty grows with the number of adapter layers: it is largest for Attn+MLP
and smallest for the half-block placements.

\paragraph{Much of the penalty is launch overhead.} The extra arithmetic of the adapters is only a few percent of
the model's, so the size of the penalty needs an explanation. Normally the CPU tells the GPU to run each operation
of the forward pass separately, a few hundred per pass, and each such \emph{launch} has a fixed cost of a few
microseconds, however small the operation. A CUDA graph records the whole sequence once and afterwards replays it
with a single launch; the GPU performs the same computations, but the cost per launch disappears. Timing the
forward pass with and without the graph therefore shows how much of it is launch overhead. On Swin-B the graph
removes 48--62\% of the extra time of the unmerged adapters; on ViT-B it removes 20--27\%. The rest is the GPU time
of the adapters' many small operations, which we did not break down further. The graph is used here only as a
diagnostic; all other latencies are measured in the normal way.

\paragraph{The penalty shrinks with batching.} The launch cost is the same whatever the batch size, so with several
images per batch it is shared among them. In batches of 8 or more, unmerged adapters cost 2--13\% of the throughput
on ViT-B and 2--17\% on Swin-B, against 8--25\% and 13--37\% at batch size 1 (Table~\ref{tab:throughput} in the
appendix). Merged models have the throughput of the dense model at every batch size.

\paragraph{On the CPU.} On the CPU, where time is mostly computation, unmerged adapters add 1--15\% (ViT-B) and
4--24\% (Swin-B) to the forward pass of one image. These values come from one process and 190 images and are rougher (the merged ViT-B models alone spread over
2\%); the CPU is the closest stand-in for an edge device that was available to us, but it is not one.

\section{What this means for deployment}

For a \textbf{single task}, the adapters should be merged before deployment. The merged model gives the same
predictions and has the size, the memory and the speed of a fully fine-tuned model. LoRA then changes nothing at
prediction time; its benefits are those of Chapter~\ref{ch:lora_ft}, during training.

Keeping the adapters \textbf{separate} is useful when several tasks share one base model on the device. Each task
then costs 2--6~MiB of storage (adapters on all blocks), and switching to it means loading and applying an adapter, which takes 2.6--7.9~ms on
ViT-B and 5.4--16.3~ms on Swin-B against 91--133~ms to load a complete model (files already in memory, so these are
lower bounds; Table~\ref{tab:merge}). The price is slower predictions, most of all for single images on the GPU,
and for MLP adapters some additional memory.

All latencies were measured on one partition of a shared server GPU, and on a server CPU. They show how the
configurations compare with each other. They are not predictions of the latency on an edge device.
